\section{Method}

\subsection{Overview}

Group-Calibrated On-Policy Distillation (GC-OPD) calibrates vanilla OPD with response-level verifier feedback while retaining its dense token-level guidance. As illustrated in Figure~\ref{fig:gcopd-overview}, GC-OPD first normalizes verifier rewards and trajectory-level OPD scores within each rollout group, then represents their difference as a signed teacher--verifier disagreement residual. Relative-advantage-based credit assignment (RACA) distributes this trajectory-level residual across tokens according to their relative OPD advantages, and the resulting correction is added to the original token-level OPD advantage.

\subsection{Group-Relative Assessments}

For each response $\mathbf{y}^{(i)}\in\mathcal{G}(\mathbf{x})$, vanilla OPD supplies a token-level OPD advantage $A^{(i)}_t$ for every response token. Equation~\ref{eq:trajectory-opd-score} aggregates these advantages into the trajectory-level OPD score $s^{(i)}$, while the verifier assigns the response a reward $R^{(i)}$. These quantities capture different properties of the response. The trajectory-level OPD score summarizes the teacher's average support relative to the frozen rollout policy, whereas the verifier reward measures verified task outcome.

Since the two signals differ in their raw scales, GC-OPD normalizes them separately within each rollout group. Let $z(\cdot)$ denote within-group z-score normalization, which subtracts the within-group mean and divides by the within-group standard deviation. We define the group-normalized verifier reward and group-normalized trajectory-level OPD score as
\begin{equation}
\begin{split}
\tilde{R}^{(i)} &=z\left (R^{(i)} \right ),\\
\tilde{s}^{(i)} &=z\left (s^{(i)} \right ).
\end{split}
\label{eq:group-relative-signals}
\end{equation}

Normalizing the two signals separately within each rollout group avoids directly comparing their raw values. For graded verifier rewards, $\tilde{R}^{(i)}$ retains the relative spacing of outcomes within the group. Binary verifier rewards, by contrast, indicate only whether each response succeeds or fails. GC-OPD operates on these group-normalized values rather than their induced rankings alone, allowing the magnitude of graded outcome differences to inform the residual.

\subsection{Teacher--Verifier Disagreement Residual}

Given the group-normalized quantities in Equation~\ref{eq:group-relative-signals}, we define the signed residual as
\begin{equation}
\rho^{(i)}=\tilde{R}^{(i)}-\tilde{s}^{(i)}.
\label{eq:residual}
\end{equation}

The sign of $\rho^{(i)}$ indicates the direction of teacher--verifier disagreement. A positive value means that the group-normalized verifier reward for response $\mathbf{y}^{(i)}$ exceeds its group-normalized OPD score, whereas a negative value indicates the opposite. Its magnitude reflects the difference between the two group-normalized assessments.

GC-OPD uses this difference instead of directly adding the group-normalized verifier reward. Direct reward addition can reinforce a relative preference already expressed by OPD. In contrast, the residual vanishes when the two group-normalized assessments coincide, while its magnitude grows with their discrepancy. The residual therefore focuses calibration on differences between the verifier and OPD assessments.

The residual also has the desired ordering under a strict disagreement between the verifier and OPD orderings. For two responses $i$ and $j$ in the same rollout group, let $\sigma_R$ and $\sigma_s$ denote the positive within-group standard deviations. If $R^{(i)}>R^{(j)}$ but $s^{(i)}<s^{(j)}$, then
\begin{equation}
\rho^{(i)}-\rho^{(j)}
=
\frac{R^{(i)}-R^{(j)}}{\sigma_R}
-
\frac{s^{(i)}-s^{(j)}}{\sigma_s}
>0.
\label{eq:pairwise-residual-direction}
\end{equation}

Thus, when the verifier prefers response $i$ but OPD prefers response $j$, the residual assigns a larger value to response $i$. Conversely, reversing the two preferences reverses the residual ordering. This verifier-consistent ordering directs the subsequent token-level correction toward the response with the higher verifier reward.

If either signal has a near-zero within-group standard deviation, the corresponding normalization is not well-defined. We therefore set $\rho^{(i)}=0$ for all responses in that group, reducing GC-OPD to vanilla OPD for the group.

\subsection{Relative-Advantage-based Credit Assignment}

The trajectory-level residual specifies the direction and magnitude of calibration for each response, but not how the correction should be distributed across tokens. RACA uses the token-level OPD advantages within each response as the allocation signal. We first normalize each token's OPD advantage relative to the response mean:
\begin{equation}
u^{(i)}_t=
\frac{A^{(i)}_t-s^{(i)}}
{\sqrt{\frac{1}{T^{(i)}}\sum_{v=1}^{T^{(i)}}
\left(A^{(i)}_v-s^{(i)}\right)^2+\epsilon}},
\end{equation}
where $\epsilon$ is a small positive constant for numerical stability. The resulting relative OPD advantage $u^{(i)}_t$ is positive for tokens whose OPD advantages exceed the response mean and negative for those below it. This normalization preserves the within-response ordering of the signed OPD advantages. Here, $u^{(i)}_t$ serves only as an allocation signal and should not be interpreted as token correctness or causal importance.

RACA maps the relative OPD advantage to a positive, bounded token credit:
\begin{equation}
c^{(i)}_t=1+\tanh\left(\frac{u^{(i)}_t}{2}\right).
\label{eq:raca}
\end{equation}

This mapping is monotonic and satisfies $c^{(i)}_t\in(0,2)$. Tokens with larger relative OPD advantages therefore receive larger credits, while $u^{(i)}_t=0$ yields unit credit $c^{(i)}_t=1$. Because the credit is always positive, it scales but does not reverse the sign of the trajectory-level residual.

The resulting GC-OPD advantage is
\begin{equation}
A'^{(i)}_t=
A^{(i)}_t+\beta c^{(i)}_t\rho^{(i)},
\label{eq:gc-opd-advantage}
\end{equation}
where $\beta\geq0$ is the residual coefficient that controls the strength of calibration. The original token-level OPD advantage remains the base term, and $\beta c^{(i)}_t\rho^{(i)}$ is the token-level residual correction. GC-OPD substitutes $A'^{(i)}_t$ for $A^{(i)}_t$ in the clipped token-level policy objective described in Section~3.1, changing the advantage construction without modifying the actor objective. Setting $\beta=0$ recovers vanilla OPD. Given the OPD advantages and verifier rewards already produced by the training pipeline, GC-OPD adds only group-level aggregation, normalization, and elementwise token transformations. It requires no additional teacher or student forward pass.
